\section{Introduction}

The United States interstate system faces two distinct categories of resilience risk. The first is structural: a corridor has no adequate alternate route when it is closed, leaving freight trapped with no viable detour. The second is environmental: a corridor is exposed to weather, flooding, or wildfire events that close it repeatedly. Paper B.1 measured the first through the B1 isolation metric --- the inverse of alternate-route availability. Paper D.1 measured the second through the D1 climate exposure metric --- composite flood, wildfire, and extreme weather vulnerability by corridor segment. This paper introduces a third concept that neither dimension captures alone: the \textit{compound resilience hole}, which occurs when a corridor scores high on both B1 and D1 simultaneously.

The compound exposure concept is not merely additive. A corridor with high B1 (no alternate) and low D1 (rarely closes) is structurally isolated but rarely tested --- its isolation is a contingency risk that may not materialize in any given year. A corridor with low B1 (alternates available) and high D1 (closes frequently) is environmentally exposed but resilient --- freight reroutes when the primary corridor closes. A compound resilience hole, by contrast, combines the worst of both: the corridor closes frequently (high D1) and when it closes, national freight has no adequate alternative (high B1). A single weather event can eliminate national freight throughput for that corridor with no rerouting option.

The operational consequences are severe. When Donner Pass closes on Interstate 80, eastbound freight must either wait --- with mean closure duration of 18 hours and maximum duration exceeding 72 hours \citep{Caltrans2023} --- or reroute via I-40 through southern California, adding 310 miles and 5.5 hours of driving. Critically, there is no interstate-standard alternate across the Sierra Nevada. US-50 and CA-88 exist but are not maintained for commercial truck traffic during winter closures. Donner closes approximately 50 times per year \citep{Caltrans2023}. The compound effect --- 50 closures, no interstate alternate, 8,000 trucks per day --- produces an estimated \$1.6 billion in annual freight disruption cost \citep{ATRI2024}.

This is the compound resilience hole in its most extreme form: a single mountain pass through which passes a dominant share of northern transcontinental freight, closing repeatedly, with no viable detour. But Donner Pass is not unique. The analysis in this paper identifies a corpus of corridors that satisfy the compound exposure condition: B1 score above 7 AND D1 score above 6, using the ROUTE v1.1 rubric \citep{ROUTE_A2}. These corridors share the property that a single weather or climate event can disrupt freight at national scale with no adequate rerouting option.

The investment implication of compound exposure differs fundamentally from single-dimension investment rationale. A corridor with high D1 alone can be addressed through infrastructure hardening: elevating roadway against flood inundation, installing snowsheds, improving drainage. A corridor with high B1 alone can be addressed through alternate route development: new alignments, upgraded parallel routes, diamond interchange zone construction \citep{ROUTE_B4}. But compound exposure corridors require \textit{both} interventions simultaneously. Addressing only D1 without B1 leaves the corridor as a single-path dependency. Addressing only B1 without D1 creates alternates that are themselves unused most of the year and underinvested relative to their strategic role. The Donner Pass freight tunnel --- a \$4 billion investment --- addresses both in a single project: it bypasses the weather-closure segment (fixing D1) and creates a second path through the Sierra Nevada (fixing B1). This dual-benefit structure gives compound exposure investments a substantially higher NPV than single-dimension investments of comparable cost.

\subsection*{Contributions}

\begin{enumerate}
\item Definition and operationalization of the compound resilience hole concept: corridors scoring B1 $>$ 7 AND D1 $>$ 6 in the ROUTE rubric
\item National corpus analysis identifying 11 corridors satisfying the compound exposure condition
\item Quantification of compound exposure costs: frequency of closure $\times$ disruption magnitude $\times$ rerouting penalty, for each priority corridor
\item Investment prioritization argument: compound exposure projects command higher NPV because a single intervention (e.g., Donner tunnel) simultaneously addresses both the structural isolation and the environmental vulnerability
\end{enumerate}
